\documentclass[12pt,a4paper,oneside]{article}
\usepackage[brazil]{babel}
\usepackage[utf8]{inputenc}
\usepackage[dvipsnames]{xcolor}
\usepackage{listings}
\usepackage{wrapfig}
\usepackage{olymp}

\setcounter{problem}{3}

\begin{document}

\begin{problem}{Copa do Mundo 2022}{copa}{2}

\begin{wrapfigure}{r}{0.25\textwidth}
\includegraphics[width=0.9\linewidth]{images/exercise_72_0.png} 
%\caption{Caption1}
%\label{fig:subim1}
\end{wrapfigure}

No final do ano passado foi disputada a Copa do Mundo 2022, no Qatar.
Esta foi a última copa com 32 seleções, divididas em 8 grupos, cada um com 4 seleções. Na próxima Copa serão 48 seleções.

Um amigo seu quis aproveitar a Copa para tentar ganhar dinheiro desenvolvendo um aplicativo para \textit{smartphone} que mostra a tabela atualizada para os usuários. Para isso, ele criou um sistema que se conecta diretamente ao site da FIFA e faz o \textit{download} dos resultados dos jogos já realizados.
Para testar o aplicativo, ele usou os dados da copa anterior, disputada na Rússia em 2018. Coincidentemente, assim como em 2022, a Suíça e a Sérvia também estavam no grupo do Brasil.

O problema é que ele se complicou na hora de fazer este sistema interpretar os resultados e calcular o número de pontos de cada seleção. Por isso, ele pediu a sua ajuda para implementar essa funcionalidade no aplicativo para ter mais sucesso em 2026, na copa que será disputada no Canadá, Estados Unidos e México.

Você deve fazer um programa que leia a lista de resultados dos jogos já realizados e calcule o número de pontos das quatro seleções que estão em um mesmo grupo.

% \Notes

% Uma seleção ganha 3 pontos para cada vitória e 1 ponto para cada empate. Derrotas não alteram sua pontuação.

\InputFile

A entrada é dividida em duas partes, separadas por uma linha em branco.

A primeira parte são os dados lidos diretamente do site da FIFA. Esta parte começa com um inteiro $N$ indicando o número de partidas já realizados. Cada uma das $N$ linhas seguintes contém o resultado de uma partida, no seguinte formato:

\begin{center}
\texttt{time1 placar1 x placar2 time2}
\end{center}

Onde:
\begin{itemize}
    \item \texttt{time1} e \texttt{time2} são os nomes das seleções (sequências de caracteres sem espaços);
    \item \texttt{placar1} é o número de gols feitos pelo \texttt{time1};
    \item \texttt{placar2} é o número de gols feitos pelo \texttt{time2};
\end{itemize}


\Notes

Para armazenar estes dados, você deve obrigatoriamente utilizar um vetor ``partida'' como definido abaixo:

\lstset{language=C++,
                basicstyle=\ttfamily,
                keywordstyle=\color{Mulberry}\ttfamily,
                stringstyle=\color{Maroon}\ttfamily,
                commentstyle=\color{YellowGreen}\ttfamily,
                morecomment=[l][\color{magenta}]{\#},
                basewidth = {.48em}
}
\begin{lstlisting}
struct TipoPartida {
    char time1[31], time2[31];
    int placar1, placar2;
};

TipoPartida partida[50];
\end{lstlisting}

A segunda parte da entrada indica quais seleções estão no grupo para o qual o aplicativo está calculando a classificação. Portanto, esta parte é composta por exatamente 4 linhas, cada uma delas indicando o nome de uma seleção. Observe que é possível que uma ou mais destas seleções ainda não tenham realizado nenhuma partida.

Restrições: $N \leq 50$; $0 \leq$ \texttt{placar1}, \texttt{placar2} $\leq 10$; todos os nomes de seleções contêm no máximo 30 caracteres e não contêm espaços.

\OutputFile

Seu programa deve gerar 4 linhas de saída, cada uma contendo o nome de uma das seleções seguida pelo número de pontos conquistados por ela até o momento. A ordem de saída deve ser a mesma de entrada, ou seja, as seleções não devem ser ordenadas de acordo com a pontuação.

% \Notes
% Preste atenção aos tipos das variáveis, pois $F(90) \approx 2^{61}$.

\Example

\begin{example}
\exmp{
3
Russia 5 x 0 ArabiaSaudita
Portugal 3 x 3 Espanha
Egito 0 x 1 Uruguai\\

Russia
ArabiaSaudita
Uruguai
Egito
}{
Russia 3
ArabiaSaudita 0
Uruguai 3
Egito 0
}%
\end{example}


\begin{example}
\exmp{
8
Russia 5 x 0 ArabiaSaudita
Egito 0 x 1 Uruguai
CostaRica 2 x 1 Servia
Argentina 1 x 1 Islandia
Brasil 0 x 0 Suica
Brasil 3 x 1 CostaRica
Servia 1 x 0 Suica
Servia 0 x 2 Brasil\\

Brasil
Suica
CostaRica
Servia
}{
Brasil 7
Suica 1
CostaRica 3
Servia 3
}%
\end{example}

\textbf{Dica}: para não ter que digitar a entrada toda vez que for testar seu programa, crie um arquivo contendo cada entrada e execute seu programa fazendo redirecionamento de entrada:

\texttt{./a.out < entrada1.txt}

Assim, toda leitura feita da entrada padrão (\texttt{scanf} e \texttt{cin}) lerá dados do arquivo \texttt{entrada1.txt}. Você não precisará digitar a entrada novamente e ela não 
aparecerá na tela, apenas o resultado do programa.
%sairá na tela, só o resultado do programa.

%Comentários sobre o exemplo, se necessário.

\end{problem}

\end{document}
